\documentclass[11pt,a4paper]{article}
\usepackage[top=1.0in, bottom=1.0in, left=1.1in, right=1.1in]{geometry}
\usepackage{graphicx}
\usepackage[sort&compress, super, numbers]{natbib}
\usepackage[export]{adjustbox}
\usepackage[ngerman]{babel}
\usepackage[utf8]{inputenc}
\usepackage[T1]{fontenc}
\renewcommand{\bibname}{References}%names reference list 
% novelty statement, a cover letter, and a title page including an authorship statement, a data statement, the number of words in abstract/main text and the number of figure/table/boxes.
% recommended and opposed reviewers, and recommended and opposed editorial board members with whom they are in conflict. 
\begin{document}

\pagenumbering{gobble}
\noindent \includegraphics[width=0.4\textwidth, right]{letterhead/ubc-logo-2018-fullsig-blue-cmyk.png}
\noindent Dear Dr. Öpik
\vspace{1.5ex}\\
\noindent Please consider our paper, ``Traits predict forest phenological responses to photoperiod more than temperature'' for publication as a full paper in \emph{New Phytologist}. 
\vspace{1.5ex}\\ 
\noindent 
Climate change is altering the environmental cues of growth globally, producing shifts in the timing of life history events or phenology, altering resource use and producing novel species interactions. But how communities respond can vary depending on individual species' traits or across spatial scales. In plants, variation in traits and environmental gradients are linked to different growth strategies, and given rise to established frameworks like the leaf or wood economic spectra \citep{Wright2004, Chave2009}. The different growth strategies within these frameworks intuitively relate to the timing of growth, but given the challenges of measuring phenological traits and their high variability, our understanding of how well phenology fits within these frameworks and whether other plant traits share the same cues is still limited.  
\vspace{1.5ex}\\  % 50 word answers
\noindent \textit{What hypotheses or questions does this work address?} 
To understand the trade-offs between traits and phenological cues require answering: how do the major budburst cues relate to traits across species and spatially? To answer this question, we combined trait measurements with large-scale budburst experiments, allowing us to estimate relationships between traits while partitioning species- and population-level differences variation.
\vspace{1.5ex}\\ 
\noindent \textit{How does this work advance our current understanding of plant science?} 
By measuring 1428 individuals from forests spanning 6$^{\circ}$ latitude and 55$^{\circ}$ longitude, our results provide the most compelling evidence of trait-phenology relationships across North America. Surprisingly, we found photoperiod to be the only cue linking phenology with broader plant growth strategies, suggesting trade-offs between traits and phenology may limit future climate responses.
\vspace{1.5ex}\\
\noindent \textit{Why is this work important and timely?} 
Our results highlight fundamental gaps in our understanding of how growth strategies and phenological cues underpin woody plant growth. The important relationships observed between traits and photoperiod---which is generally considered a weak cue of budburst---will improve our ability to identify how species growth will change with climate change and future vulnerability.
\vspace{1.5ex}\\
\noindent All authors contributed to this work and approve this version for submission. The manuscript is 4083 words with a 200 word summary, and 4 figures and is not under consideration elsewhere. We hope you find it suitable for publication in \emph{New Phytologist}, and look forward to hearing from you. 
%\noindent We recommend the following reviewers: Dr. Meredith Zettlemoyer, Dr. Mason Heberling, Dr. Jason Fridley, Dr. Rong Yu, and Dr. Ameila Caffara. 
\vspace{1.5ex}\\
\noindent Sincerely, \\
\includegraphics[scale=.4]{letterhead/shot.png} \\ 
\noindent Deirdre Loughnan\\
\noindent Sentinels of Change Postdoctoral Fellow\\ %emw19Dec: Nice!
\noindent Hakai Institute $|$ Department of Zoology\\
\noindent University of British Columbia
\newpage
\vspace{-5ex}

\begingroup
\renewcommand{\section}[2]{}
\bibliographystyle{refs/bibstyles/nature.bst}% 
\bibliography{refs/biblio.bib}
\endgroup
\newpage


\end{document}

%\begingroup
%\renewcommand{\section}[2]{}%
%\renewcommand{\chapter}[2]{}% for other classes
%\bibliography{mybibfile}
%\endgroup